\section{Executed Study: Detectability Across Evidence Tiers (RQ2--RQ4)}
\label{sec:study}

We now measure the taxonomy's detectability structure. The study is executed
and deterministically reproducible from the packaged source
($\approx$380 lines of dependency-free Python; the results table, including
its header, is generated programmatically by the same run). Its
epistemic scope is stated plainly: it is a \emph{closed-world evaluation of
detector semantics}---the world, its scenarios, and the detectors are all
models---so it establishes what each evidence tier can and cannot decide
\emph{in principle}, with event-time latencies. It does not measure any real
cluster; wall-clock behavior on real infrastructure is the laboratory's job
(\cref{sec:lab}).

\subsection{Modeled World and Scenarios}
The world models one deployment with the selected admitted basis of
\cref{sec:model-states}---including recorded environment assumptions
$\sigma_B$---and streams for desired state (Git), observed cluster state,
registry contents, policy versions \emph{with evaluable requirement sets}
(so that T1 implements \cref{def:poldrift}'s admission re-evaluation, not a
version-string comparison), authorization records as a set with subject
digests, covered revisions, windows, and revocation (so that legitimate
successors are modeled), and two unmanaged environment properties. Each run
is 400 events; drift is injected at event 150. Twelve scenarios follow the
taxonomy and the constructed examples of \cref{fig:example}:

\begin{description}[leftmargin=2.6em,itemsep=1pt,topsep=2pt]
\item[S1] manual in-cluster change (observed field diverges from desired);
\item[S2] expired emergency authorization (exception granted with a 40-event
window; nothing removes it at expiry);
\item[S3] policy version change post-deploy ($\pi_7 \to \pi_8$, which the
running configuration violates);
\item[S4] artifact substitution (same tag re-pointed to a new digest,
rolled onto nodes; manifests unchanged);
\item[S5] IAM permission expansion out of band (unmanaged);
\item[S6] Git rollback without new approval (reconciler converges to the
reverted content; approval binds the newer revision);
\item[S7] out-of-band cloud console modification (unmanaged);
\item[S8] approval mismatch (deployed digest is not any approval's subject);
\item[S9] live-status deletion for a continuing approval: immutable
$A_B^{\mathrm{proof}}$ survives, but the authenticated source required to
establish continuing validity is lost; authorization and intent become
undecidable---the missing status is neither a revocation nor proof that approval never existed
(evidence drift---\cref{def:evddrift}---exercised directly);
\item[S10] policy supersession plus an expired exception
($\{\text{policy},\text{authorization}\}$);
\item[S11] artifact substitution plus environment change
($\{\text{authorization},\text{environment}\}$);
\item[S12] rollback plus deletion of its approval basis: the underlying
intent/authorization relations become undecidable, so the honest observable
output is evidence drift rather than a guessed polar set;
\end{description}
plus \textbf{S0}, a no-drift control. Churn is of two kinds, both present in
every stream including the control: \emph{runtime churn} (30\% of events:
autoscaler replica changes, rolling-restart hash changes) and \emph{benign
governance churn} (2\% of events: policy revisions the running state
satisfies, approved updates that legitimately advance revision, digest, and
approval together, and hygienic exceptions removed at expiry). The governance
churn exists so that the governance tiers' false-alarm rates are earned
against realistic governance noise rather than measured against a silent
plane.

\subsection{Detectors and Counterfactual Scoring}
Six detector tiers implement \cref{sec:architecture}: T0 (naive
configuration comparison, including runtime-owned fields), T0n (normalized),
and cumulative T1--T4, each a pure function of its tier's visible streams, so
tier capabilities are enforced by construction. Ground truth is a
\emph{class set} per scenario (components are a vector and may drift
together, \cref{def:govdrift}). The detector exports the complete class set
decidable at its tier; a priority-ordered first verdict is derived only as an
operational interface. We score exact-set agreement, subset-only partial
diagnosis, and Hamming loss rather than accepting any one member of the ground
truth. Missing approval records set the relevant entries of $\ObsMask$ to zero
and produce the explicit undecidable/evidence verdict, never a guessed polar
one. S6, S10, and S11 test substantive composition; S12 tests composition under
evidence loss.

Scoring required care, and the naive arm is why. A detector that alarms
constantly ``detects'' every scenario by coincidence; volume is not
information. We therefore score \emph{counterfactually}: for each (scenario,
tier, seed), the detector runs on the drift stream and on a paired,
same-seed, no-drift control stream with identical churn; the drift is
\emph{detected} iff the two alarm sequences differ, detection time is the
first differing event, and classification is the class set emitted there. Under
this rule a constant alarmer carries zero information and scores zero. False
alarms are counted on the control stream. Twenty seeds per cell; onset for S2
is the expiry instant (the drift is the exception \emph{outliving} its
window, not its grant).

\subsection{Results}
\begin{table*}[t]
\centering
\caption{Detectability Matrix, Generated Programmatically from the Run
Outputs (20 Seeds per Cell). \cmark{} = Detected with Exact-Set Agreement in
All Seeds; $\triangle$ = Detected but Only a Correct Subset Is Decidable at
That Tier (Superscript: Mean Detection Latency in Events When $>$0);
\xmark{} = Not Detected (No Information Under Counterfactual Scoring). FP =
False Alarms per 400-Event Control Run.}
\label{tab:matrix}
\footnotesize
\resizebox{0.96\textwidth}{!}{\begin{tabular}{@{}lcccccc@{}}
\toprule
\textbf{Scenario} & \textbf{T0} & \textbf{T0n} & \textbf{T1} & \textbf{T2} & \textbf{T3} & \textbf{T4} \\
 & {\scriptsize naive cfg} & {\scriptsize norm.\ cfg} & {\scriptsize $+$policy} & {\scriptsize $+$authz} & {\scriptsize $+$lineage} & {\scriptsize $+$env} \\
\midrule
S0 control (churn only) & FP: 285.4 & FP: 0.0 & FP: 0.0 & FP: 0.0 & FP: 0.0 & FP: 0.0 \\
S1 & \cmark$^{+22.6}$ & \cmark & \cmark & \cmark & \cmark & \cmark \\
S2 & \xmark & \xmark & \xmark & \cmark & \cmark & \cmark \\
S3 & \xmark & \xmark & \cmark & \cmark & \cmark & \cmark \\
S4 & \xmark & \xmark & \xmark & \xmark & \cmark & \cmark \\
S5 & \xmark & \xmark & \xmark & \xmark & \xmark & \cmark \\
S6 & \xmark & \xmark & \xmark & \cmark & \cmark & \cmark \\
S7 & \xmark & \xmark & \xmark & \xmark & \xmark & \cmark \\
S8 & \xmark & \xmark & \xmark & \xmark & \cmark & \cmark \\
S9 & \xmark & \xmark & \xmark & \cmark & \cmark & \cmark \\
\bottomrule
\end{tabular}
}
\par\vspace{2pt}
\parbox{0.96\textwidth}{\scriptsize S1 manual change; S2 expired exception; S3
policy supersession; S4 artifact substitution; S5 IAM expansion; S6
unapproved rollback; S7 out-of-band cloud change; S8 approval mismatch; S9
approval-record deletion; S10 policy+exception; S11 artifact+environment;
S12 rollback+missing approval. Tiers are cumulative; T0n normalizes
runtime-owned fields; latencies use expiry-inclusive semantics.}
\end{table*}

\Cref{tab:matrix} is the study's central artifact. Four findings, stated in
the register the design supports.

\begin{finding}[The staircase: implementation validation of tier
dependencies]
\label{fnd:staircase}
Detection forms a strict staircase realizing
\cref{prop:tiers,prop:tier-minimality} cell for cell: T0n detects S1 and
nothing else---one scenario of twelve; T1 adds policy drift (S3) and the
policy component of S10, and---because it re-evaluates policy rather than
comparing versions---raises no alarm on satisfied policy revisions in the
churn. T2 adds the validity- and coverage-dependent scenarios S2, S6, and
S10, plus the evidence-decay scenarios S9 and S12 (undecidable, not polar).
T3 adds the digest-level scenarios S4, S8, and S11 and completes S6's
two-component diagnosis; T4 adds S5/S7 and completes S11. The triangle cells
show why first-label accuracy was insufficient: a lower tier can report a
correct subset without deciding the full vector. Every reported component is
correct in every seed, and exact-set agreement appears when the tier carries
all required streams. We state this finding's epistemic value precisely: the world and
detectors are co-designed, so the exactness is expected, and the staircase is
\emph{implementation validation} of the model's dependency structure plus a
noise-behavior measurement---not field evidence for the taxonomy, a role we
do not claim for it. Its concrete content for RQ4 stands: in the reference
semantics, configuration comparison decides only its own substantive
component; the five-class substantive vector plus epistemic mask requires the
joined tiers. All six umbrella classes are exercised.
\end{finding}

\begin{finding}[A deliberately naive comparator validates the scoring rule]
\label{fnd:naive}
The unnormalized diagnostic comparator T0 alarms on 285.4 of 400 control events
(71.4\% in our churn model)---and the rate is nearly flat across runtime
churn rates of 10--50\% (287, 285, 300 alarms per run), because alarms
track the \emph{dwell} of divergent runtime-owned state, not the event
rate. Under counterfactual scoring T0 still detects S1, but with mean
latency 22.6 events versus 0 for T0n: its information is buried in its own
noise. For the other eleven scenarios it carries zero information while
alarming continuously. T0 is not an industry baseline---production
reconcilers normalize---and no comparative effectiveness claim rests on it;
its sole role is to demonstrate that the counterfactual scoring rule does not
reward alarm volume. Conversely, the governance tiers' zero false
alarms are now \emph{earned}: the control contains satisfied policy
revisions (which a version-comparison T1 would have flagged), legitimate
approved updates (which a pinned-revision intent check would have flagged),
and hygienic exceptions---and the definition-faithful detectors correctly
ignore all of them.
\end{finding}

\begin{finding}[Governance drift is event-detectable, with one designed
exception]
\label{fnd:latency}
At the correct tier, every scenario is detected with latency 0 events
(including S2 under expiry-inclusive semantics) and none requires scanning: the deciding facts arrive as
events on governance streams (a policy revision, a clock crossing, a registry
event, an inventory delta). RQ3's answer in the closed world is therefore:
visibility is bounded by evaluation cadence on the governance plane, not by
any intrinsic latency of the phenomenon---which sharpens what the laboratory
must measure (wall-clock cadence effects, \cref{sec:lab}).
\end{finding}

\begin{finding}[Reconciliation health is compatible with total governance
failure]
\label{fnd:blind}
In eleven of twelve scenarios, post-convergence normalized configuration
comparison reports clean---while the
deployment runs under superseded policy, expired authorization, unapproved
content, or altered environment. This is \cref{prop:impossible} made
operational, and it quantifies the blind spot claimed in
\cref{sec:background}: a \textsc{Synced} reconciler is evidence about
configuration, and about configuration only.
\end{finding}

\subsection{Independent Composite Baseline: Signals Without the Join}
The staircase establishes dependency realization but does not isolate the
value of the admitted-basis join. We therefore execute a second baseline that
unions five plane-local mechanisms: normalized desired/observed comparison,
current-policy evaluation, exception-expiry checking, digest membership in a
tool-local attestation trust set, and current environment posture against a
benchmark. It receives the same event stream and benign churn as T4 but has
no $B_x(t)$, snapshot selection, intent history, or evidence-state semantics.
This is a stronger baseline than T0 and represents composition of existing
signals without claiming to reproduce a particular commercial product.

\begin{table*}[t]
\centering
\caption{Composite Plane-Local Baseline Versus the Admitted-Basis Join.
Exact Columns Are Percentages Over 20 Seeds; Empty Local Output on S9 Means
the Deleted Governance Stream Is Outside Its Vocabulary.}
\label{tab:join-baseline}
\footnotesize
\begin{tabular}{@{}lccc rr@{}}
\toprule
Scenario & Expected & Local union & Joined & Local exact & Joined exact \\
\midrule
S1 & $\{config.\}$ & $\{config.\}$ & $\{config.\}$ & 100 & 100 \\
S2 & $\{auth.\}$ & $\{auth.\}$ & $\{auth.\}$ & 100 & 100 \\
S3 & $\{policy\}$ & $\{policy\}$ & $\{policy\}$ & 100 & 100 \\
S4 & $\{auth.\}$ & $\{auth.\}$ & $\{auth.\}$ & 100 & 100 \\
S5 & $\{env.\}$ & $\{env.\}$ & $\{env.\}$ & 100 & 100 \\
S6 & $\{intent,auth.\}$ & $\{auth.\}$ & $\{intent,auth.\}$ & 0 & 100 \\
S7 & $\{env.\}$ & $\{env.\}$ & $\{env.\}$ & 100 & 100 \\
S8 & $\{auth.\}$ & $\{auth.\}$ & $\{auth.\}$ & 100 & 100 \\
S9 & $\{evidence\}$ & $\{\}$ & $\{evidence\}$ & 0 & 100 \\
S10 & $\{policy,auth.\}$ & $\{policy,auth.\}$ & $\{policy,auth.\}$ & 100 & 100 \\
S11 & $\{auth.,env.\}$ & $\{auth.,env.\}$ & $\{auth.,env.\}$ & 100 & 100 \\
S12 & $\{evidence\}$ & $\{auth.\}$ & $\{evidence\}$ & 0 & 100 \\
\bottomrule
\end{tabular}

\end{table*}

The local union exactly identifies 9/12 scenario sets (180/240
scenario-seed units, 75\%); joined T4 identifies 12/12 (240/240, 100\%).
Both remain quiet on the paired controls, so the difference is diagnosis,
not alarm volume. The local union misses S6's intent component, cannot see
S9's evidence loss, and mislabels S12 as an authorization finding rather than
withholding a polar verdict. Four semantic probes sharpen causality: identical
current-policy failure with versus without an admitted basis produces the
same local ``policy'' alert, while the join returns policy versus evidence;
missing live status yields evidence for a continuing instrument but remains
consistent for a prospective one-shot instrument with retained proof. The
incremental value is therefore the history- and mode-aware relation that
selects both class and remediation family. As with the staircase, this
closed-world comparison validates semantics, not field effectiveness.

\subsection{Threats to Validity}
\emph{Construct}: detectors and world share a modeling vocabulary; the tiers'
\emph{exact} 0/1 staircase is a property of the closed world, where streams
are noiseless and semantics are implemented as defined. Real streams are
late, lossy, and ambiguous; the staircase's \emph{shape} (which class needs
which evidence) is the transferable claim---it follows from input
dependencies (\cref{prop:tiers})---while the crispness is not.
\emph{Internal}: single-deployment world and one injection time per run. Three
compound scenarios exercise set output, but do not span the combinatorial
space. The first-verdict order is one operational policy among several;
exact-set scoring is invariant to it. Scenario S6's transient configuration signal
(before the reconciler converges) is not credited to T0n because the
comparison evaluates post-convergence state in our event ordering; on real
clusters a fast comparator might glimpse it---a cadence question for the lab.
\emph{External}: the repeated laboratory (\cref{sec:lab}) shows that the nine
scenario classes can be induced and usually observed through real
Kubernetes/GitOps interfaces on one pinned stack. It does not transfer the
closed-world rates quantitatively or estimate natural occurrence, field
prevalence, stream-loss reliability, multi-cluster behavior, or performance
under load. Its two cadence misses are direct evidence against assuming that
polling necessarily observes transient configuration drift.
